% GENERATED by tools/xref.py from the .aux files of the other volumes; do not edit.
\xrlabel{chap:1}{{1}{}{Light dark matter}{Doc-Start}{}}
\xrlabel{chap:1@cref}{{[chapter][1][]1}{[1][1][]1}}
\xrlabel{sec:1.1}{{1.1}{}{The master equation of direct detection}{Doc-Start}{}}
\xrlabel{sec:1.1@cref}{{[section][1][1]1.1}{[1][1][]1}}
\xrlabel{sec:1.2}{{1.2}{}{The local halo}{Doc-Start}{}}
\xrlabel{sec:1.2@cref}{{[section][2][1]1.2}{[1][1][]1}}
\xrlabel{tab:1-shm}{{1.1}{}{Standard-halo-model parameters adopted in this book, with their sources. Vectors are given in galactic axes (toward the Galactic center, along Galactic rotation, toward the north Galactic pole). The last two rows follow from the others. The reference Earth speed is the adopted magnitude of the solar velocity $(0, 238, 0) + (11.1, 12.2, 7.3)$\,\si {km.s^{-1}}. The largest laboratory speed adds it to the escape speed}{Doc-Start}{}}
\xrlabel{tab:1-shm@cref}{{[table][1][1]1.1}{[1][1][]1}}
\xrlabel{eq:1.2-vsun}{{1.5}{}{Density and velocity distribution}{Doc-Start}{}}
\xrlabel{eq:1.2-vsun@cref}{{[equation][5][1]1.5}{[1][1][]1}}
\xrlabel{eq:1.2-flab}{{1.8}{}{The laboratory frame}{Doc-Start}{}}
\xrlabel{eq:1.2-flab@cref}{{[equation][8][1]1.8}{[1][1][]1}}
\xrlabel{eq:1.2-vmax}{{1.9}{}{The laboratory frame}{Doc-Start}{}}
\xrlabel{eq:1.2-vmax@cref}{{[equation][9][1]1.9}{[1][1][]1}}
\xrlabel{eq:1.2-rotation}{{1.14}{}{Daily modulation}{Doc-Start}{}}
\xrlabel{eq:1.2-rotation@cref}{{[equation][14][1]1.14}{[1][1][]1}}
\xrlabel{sec:1.3}{{1.3}{}{Kinematics of slow dark matter}{Doc-Start}{}}
\xrlabel{sec:1.3@cref}{{[section][3][1]1.3}{[1][1][]1}}
\xrlabel{eq:1.3-omegaq}{{1.21}{}{Energy transfer and minimum velocity}{Doc-Start}{}}
\xrlabel{eq:1.3-omegaq@cref}{{[equation][21][1]1.21}{[1][1][]1}}
\xrlabel{eq:1.3-vmin}{{1.22}{}{Energy transfer and minimum velocity}{Doc-Start}{}}
\xrlabel{eq:1.3-vmin@cref}{{[equation][22][1]1.22}{[1][1][]1}}
\xrlabel{eq:1.3-mmin}{{1.24}{}{Threshold mass}{Doc-Start}{}}
\xrlabel{eq:1.3-mmin@cref}{{[equation][24][1]1.24}{[1][1][]1}}
\xrlabel{eq:1.3-eta}{{1.27}{}{The speed integral}{Doc-Start}{}}
\xrlabel{eq:1.3-eta@cref}{{[equation][27][1]1.27}{[1][1][]1}}
\xrlabel{sec:1.3-edge}{{1.3.5}{}{The kinematic edge}{Doc-Start}{}}
\xrlabel{sec:1.3-edge@cref}{{[subsection][5][1,3]1.3.5}{[1][1][]1}}
\xrlabel{sec:1.4}{{1.4}{}{Light dark-matter models}{Doc-Start}{}}
\xrlabel{sec:1.4@cref}{{[section][4][1]1.4}{[1][1][]1}}
\xrlabel{eq:1.4-yukawa}{{1.30}{}{Mediators and the Yukawa potential}{Doc-Start}{}}
\xrlabel{eq:1.4-yukawa@cref}{{[equation][30][1]1.30}{[1][1][]1}}
\xrlabel{eq:1.4-Vq}{{1.31}{}{Mediators and the Yukawa potential}{Doc-Start}{}}
\xrlabel{eq:1.4-Vq@cref}{{[equation][31][1]1.31}{[1][1][]1}}
\xrlabel{eq:1.4-defs}{{1.32}{}{Reference cross-section and form factor}{Doc-Start}{}}
\xrlabel{eq:1.4-defs@cref}{{[equation][32][1]1.32}{[1][1][]1}}
\xrlabel{eq:1.4-Vsq}{{1.33}{}{Reference cross-section and form factor}{Doc-Start}{}}
\xrlabel{eq:1.4-Vsq@cref}{{[equation][33][1]1.33}{[1][1][]1}}
\xrlabel{eq:1.4-sigmaN}{{1.36}{}{Scattering on nuclei}{Doc-Start}{}}
\xrlabel{eq:1.4-sigmaN@cref}{{[equation][36][1]1.36}{[1][1][]1}}
\xrlabel{sec:1.5}{{1.5}{}{Detector signals}{Doc-Start}{}}
\xrlabel{sec:1.5@cref}{{[section][5][1]1.5}{[1][1][]1}}
\xrlabel{eq:1.5-rate}{{1.46}{}{Event rate}{Doc-Start}{}}
\xrlabel{eq:1.5-rate@cref}{{[equation][46][1]1.46}{[1][1][]1}}
\xrlabel{sec:2.2}{{2.2}{}{Born approximation and form factors}{Doc-Start}{}}
\xrlabel{sec:2.2@cref}{{[section][2][2]2.2}{[1][1][]1}}
\xrlabel{eq:2.2-dsdq}{{2.15}{}{The Born cross-section}{Doc-Start}{}}
\xrlabel{eq:2.2-dsdq@cref}{{[equation][15][2]2.15}{[1][1][]1}}
\xrlabel{eq:2.2-fion}{{2.22}{}{The hydrogen atom}{Doc-Start}{}}
\xrlabel{eq:2.2-fion@cref}{{[equation][22][2]2.22}{[1][1][]1}}
\xrlabel{eq:2.2-fion-pw}{{2.24}{}{The hydrogen atom}{Doc-Start}{}}
\xrlabel{eq:2.2-fion-pw@cref}{{[equation][24][2]2.24}{[1][1][]1}}
\xrlabel{eq:2.2-fermi}{{2.25}{}{The hydrogen atom}{Doc-Start}{}}
\xrlabel{eq:2.2-fermi@cref}{{[equation][25][2]2.25}{[1][1][]1}}
\xrlabel{eq:2.2-nucleus}{{2.27}{}{Nuclei and the Helm form factor}{Doc-Start}{}}
\xrlabel{eq:2.2-nucleus@cref}{{[equation][27][2]2.27}{[1][1][]1}}
\xrlabel{sec:2.3}{{2.3}{}{Second quantization}{Doc-Start}{}}
\xrlabel{sec:2.3@cref}{{[section][3][2]2.3}{[1][1][]1}}
\xrlabel{eq:2.3-slater}{{2.29}{}{Identical particles}{Doc-Start}{}}
\xrlabel{eq:2.3-slater@cref}{{[equation][29][2]2.29}{[1][1][]1}}
\xrlabel{eq:2.3-Hee-density}{{2.44}{}{Coulomb interaction in momentum space}{Doc-Start}{}}
\xrlabel{eq:2.3-Hee-density@cref}{{[equation][44][2]2.44}{[1][1][]1}}
\xrlabel{we:2-jellium}{{2.5}{}{Coulomb Hamiltonian of a solid}{Doc-Start}{}}
\xrlabel{we:2-jellium@cref}{{[workedexample][5][2]2.5}{[1][1][]1}}
\xrlabel{eq:2.3-jellium}{{2.47}{}{Coulomb Hamiltonian of a solid}{Doc-Start}{}}
\xrlabel{eq:2.3-jellium@cref}{{[equation][47][2]2.47}{[1][1][]1}}
\xrlabel{eq:2.3-HC}{{2.48}{}{Coulomb Hamiltonian of a solid}{Doc-Start}{}}
\xrlabel{eq:2.3-HC@cref}{{[equation][48][2]2.48}{[1][1][]1}}
\xrlabel{sec:2.4}{{2.4}{}{Second-order processes}{Doc-Start}{}}
\xrlabel{sec:2.4@cref}{{[section][4][2]2.4}{[1][1][]1}}
\xrlabel{sec:2.5}{{2.5}{}{Thermal states and detailed balance}{Doc-Start}{}}
\xrlabel{sec:2.5@cref}{{[section][5][2]2.5}{[1][1][]1}}
\xrlabel{eq:2.5-golden}{{2.67}{}{Thermal golden rule}{Doc-Start}{}}
\xrlabel{eq:2.5-golden@cref}{{[equation][67][2]2.67}{[1][1][]1}}
\xrlabel{eq:2.5-fermi}{{2.71}{}{Free electron gas at zero temperature}{Doc-Start}{}}
\xrlabel{eq:2.5-fermi@cref}{{[equation][71][2]2.71}{[1][1][]1}}
\xrlabel{sec:3.1}{{3.1}{}{Lattices and reciprocal space}{Doc-Start}{}}
\xrlabel{sec:3.1@cref}{{[section][1][3]3.1}{[1][1][]1}}
\xrlabel{eq:3.1-kgrid}{{3.13}{}{Crystal momentum and the lattice sum}{Doc-Start}{}}
\xrlabel{eq:3.1-kgrid@cref}{{[equation][13][3]3.13}{[1][1][]1}}
\xrlabel{eq:3.1-latticesum}{{3.14}{}{Crystal momentum and the lattice sum}{Doc-Start}{}}
\xrlabel{eq:3.1-latticesum@cref}{{[equation][14][3]3.14}{[1][1][]1}}
\xrlabel{we:3-diamond}{{3.2}{}{The basis structure factor}{Doc-Start}{}}
\xrlabel{we:3-diamond@cref}{{[workedexample][2][3]3.2}{[1][1][]1}}
\xrlabel{sec:3.2}{{3.2}{}{Bloch states and band structure}{Doc-Start}{}}
\xrlabel{sec:3.2@cref}{{[section][2][3]3.2}{[1][1][]1}}
\xrlabel{eq:3.2-He}{{3.19}{}{Bloch's theorem}{Doc-Start}{}}
\xrlabel{eq:3.2-He@cref}{{[equation][19][3]3.19}{[1][1][]1}}
\xrlabel{eq:3.2-bloch}{{3.20}{}{Bloch's theorem}{Doc-Start}{}}
\xrlabel{eq:3.2-bloch@cref}{{[equation][20][3]3.20}{[1][1][]1}}
\xrlabel{eq:3.2-Hk}{{3.21}{}{Bloch's theorem}{Doc-Start}{}}
\xrlabel{eq:3.2-Hk@cref}{{[equation][21][3]3.21}{[1][1][]1}}
\xrlabel{eq:3.2-pw}{{3.22}{}{Bloch's theorem}{Doc-Start}{}}
\xrlabel{eq:3.2-pw@cref}{{[equation][22][3]3.22}{[1][1][]1}}
\xrlabel{eq:3.2-central}{{3.23}{}{Bloch's theorem}{Doc-Start}{}}
\xrlabel{eq:3.2-central@cref}{{[equation][23][3]3.23}{[1][1][]1}}
\xrlabel{eq:3.2-blochsum}{{3.28}{}{Tight binding and the \texorpdfstring {$sp^3$}{sp3} bond}{Doc-Start}{}}
\xrlabel{eq:3.2-blochsum@cref}{{[equation][28][3]3.28}{[1][1][]1}}
\xrlabel{eq:3.2-HS}{{3.29}{}{Tight binding and the \texorpdfstring {$sp^3$}{sp3} bond}{Doc-Start}{}}
\xrlabel{eq:3.2-HS@cref}{{[equation][29][3]3.29}{[1][1][]1}}
\xrlabel{eq:3.2-LCAO}{{3.30}{}{Tight binding and the \texorpdfstring {$sp^3$}{sp3} bond}{Doc-Start}{}}
\xrlabel{eq:3.2-LCAO@cref}{{[equation][30][3]3.30}{[1][1][]1}}
\xrlabel{eq:3.2-rho}{{3.34}{}{The density vertex}{Doc-Start}{}}
\xrlabel{eq:3.2-rho@cref}{{[equation][34][3]3.34}{[1][1][]1}}
\xrlabel{eq:3.2-rhopw}{{3.35}{}{The density vertex}{Doc-Start}{}}
\xrlabel{eq:3.2-rhopw@cref}{{[equation][35][3]3.35}{[1][1][]1}}
\xrlabel{eq:3.2-r}{{3.39}{}{Velocity matrix elements}{Doc-Start}{}}
\xrlabel{eq:3.2-r@cref}{{[equation][39][3]3.39}{[1][1][]1}}
\xrlabel{eq:3.2-dHk}{{3.40}{}{Velocity matrix elements}{Doc-Start}{}}
\xrlabel{eq:3.2-dHk@cref}{{[equation][40][3]3.40}{[1][1][]1}}
\xrlabel{eq:3.2-vg}{{3.41}{}{Velocity matrix elements}{Doc-Start}{}}
\xrlabel{eq:3.2-vg@cref}{{[equation][41][3]3.41}{[1][1][]1}}
\xrlabel{eq:3.2-dudk}{{3.42}{}{Velocity matrix elements}{Doc-Start}{}}
\xrlabel{eq:3.2-dudk@cref}{{[equation][42][3]3.42}{[1][1][]1}}
\xrlabel{eq:3.2-optical}{{3.43}{}{Velocity matrix elements}{Doc-Start}{}}
\xrlabel{eq:3.2-optical@cref}{{[equation][43][3]3.43}{[1][1][]1}}
\xrlabel{sec:3.3}{{3.3}{}{Silicon, germanium and gallium arsenide}{Doc-Start}{}}
\xrlabel{sec:3.3@cref}{{[section][3][3]3.3}{[1][1][]1}}
\xrlabel{fig:3-si-bands}{{3.7}{}{Silicon. (a)~Band structure along $\mathrm {L}$--$\mathrm {\Gamma }$--$\mathrm {X}$--$\mathrm {W}$--$\mathrm {K}$--$\mathrm {\Gamma }$ (points of \cref {sec:3.1}), valence bands black and conduction bands blue, energies from the valence-band maximum. The bands are the Kohn--Sham bands of the plane-wave calculation of \cref {app:E.1} at the measured lattice constant: generalized-gradient functional of Perdew, Burke and Ernzerhof (PBE), SG15 norm-conserving pseudopotential, cutoff \SI {60}{Ry}, $12^3$ grid for the density, $a = \SI {5.431}{\angstrom }$, no spin--orbit coupling~\cite {giannozzi2017advanced,hamann2013optimized,schlipf2015optimization}. They put the conduction-band minimum at $k_{\mathrm {c}} = 0.843\,(2\pi /a)$ (measured $0.85$) and the indirect gap $\Eg $ at \SI {0.557}{eV} (measured \SI {1.17}{eV}), the gap problem of density-functional theory. The vertical gap at $\mathrm {\Gamma }$, marked $E_{\mathrm {dir}}$, is \SI {2.56}{eV} (measured \SI {3.4}{eV}) and is the smallest vertical gap in the zone; the one at $\mathrm {L}$ is \SI {2.70}{eV} (measured \SI {3.45}{eV}). (b)~Minimum pair energy $E_{\mathrm {min}}(\vect {q})$, \cref {eq:3.3-Emin}, for $\vect {q}$ along [100], [110] and [111], from the same calculation on a grid of $48^3$ crystal momenta, with the conduction bands shifted rigidly up by \SI {0.61}{eV} (a scissor, \cref {sec:3.4}) so that the indirect gap equals the measured \SI {1.17}{eV} (dash-dotted line); at $\vect {q} = 0$ the shifted bands give $E_{\mathrm {min}} = \SI {3.17}{eV}$. Dashed: the largest energy that dark matter of \SI {1}{MeV} can deposit, $qv_{\mathrm {max}} - q^2/2\mchi $; dotted: the bound $qv_{\mathrm {max}}$ for any mass, with $v_{\mathrm {max}} = 2.650\times 10^{-3}$. The vertical line marks $|\vect {k}_{\mathrm {c}}| = \SI {1.94}{keV}$. Along [100], $E_{\mathrm {min}}$ returns to its value at $\vect {q} = 0$ at $|\vect {G}_{200}| = \SI {4.57}{keV}$}{Doc-Start}{}}
\xrlabel{fig:3-si-bands@cref}{{[figure][7][3]3.7}{[1][1][]1}}
\xrlabel{fig:3-ge-gaas}{{3.8}{}{Band structures of (a)~germanium and (b)~GaAs along the path of \cref {fig:3-si-bands}(a). Energies from the valence-band maximum; red dots mark the lowest conduction states away from $\mathrm {\Gamma }$. (a)~Local empirical pseudopotential of Cohen and Bergstresser~\cite {cohen1966band} at $a = \SI {5.66}{\angstrom }$, the lattice constant of their fit, without spin--orbit coupling, with the conduction bands shifted down by \SI {0.21}{eV} so that the $\mathrm {L}$ minimum sits at the measured \SI {0.744}{eV}. The crystal is indirect, with the $s$-like state $\mathrm {\Gamma }_{2'c}$ (square) at \SI {1.01}{eV} (measured \SI {0.89}{eV}) and the $\Delta $ minimum, at $0.83\,(2\pi /a)$ along $\mathrm {\Gamma }$--$\mathrm {X}$, at \SI {0.81}{eV}. Plain PBE with spin--orbit coupling at the measured lattice constant closes the germanium gap, with $\mathrm {\Gamma }_{7c}$ \SI {0.05}{eV} below the valence maximum. (b)~PBE plane-wave calculation of \cref {fig:3-si-bands} (SG15 pseudopotentials with the Ga $3d$ shell in the valence, \SI {60}{Ry}, $12^3$ grid, no spin--orbit coupling) at the room-temperature lattice constant $a = \SI {5.653}{\angstrom }$~\cite {giannozzi2017advanced,hamann2013optimized,schlipf2015optimization}. The gap is direct, as measured, but PBE underestimates it: \SI {0.51}{eV} against the measured \SI {1.52}{eV} at low temperature, with the $\mathrm {L}$ and $\mathrm {X}$ states at 1.00 and \SI {1.46}{eV} (measured 1.815 and \SI {1.981}{eV})}{Doc-Start}{}}
\xrlabel{fig:3-ge-gaas@cref}{{[figure][8][3]3.8}{[1][1][]1}}
\xrlabel{tab:3-materials}{{3.1}{}{Band parameters of silicon, germanium and GaAs. Lattice constant $a$ at room temperature: silicon~\cite [lattice parameter of silicon]{mohr2025codata}, germanium~\cite [Sect.~II\,A]{rideau2006strained}, GaAs~\cite {vurgaftman2001band}. Gaps at low temperature: silicon~\cite [Eq.~(11)]{ramanathan2020ionization}, germanium~\cite [Table~VI]{rideau2006strained}, GaAs~\cite [Table~I]{lucero2012improved}; at \SI {300}{K}~\cite [Table~I]{lucero2012improved}. Minimum direct gaps $E_{\mathrm {dir}}$ and spin--orbit splittings $\Delta _{\mathrm {SO}}$ of silicon and germanium~\cite [Table~II]{rideau2006strained}. Conduction valleys: position and number; masses along ($m^*_\parallel $) and across ($m^*_\perp $) the valley axis, silicon~\cite {hensel1965cyclotron} as listed in Ref.~\cite [Table~VII]{rideau2006strained}, germanium~\cite [Table~VII]{rideau2006strained}. Heavy- and light-hole masses $m_{\mathrm {hh}}$, $m_{\mathrm {lh}}$ along [100] and [111] from \cref {eq:3.3-holes} with the Luttinger parameters $(\gamma ^{\mathrm {Lu}}_1, \gamma ^{\mathrm {Lu}}_2, \gamma ^{\mathrm {Lu}}_3) = (4.285, 0.339, 1.446)$ for silicon and $(13.25, 4.20, 5.56)$ for germanium~\cite [Table~VII]{rideau2006strained}, and $(6.98, 2.06, 2.93)$ for GaAs. All GaAs entries not credited otherwise are from Ref.~\cite {vurgaftman2001band}. Masses in units of $m_e$}{Doc-Start}{}}
\xrlabel{tab:3-materials@cref}{{[table][1][3]3.1}{[1][1][]1}}
\xrlabel{eq:3.3-kane2}{{3.52}{}{Effective masses from \texorpdfstring {$\vect {k}\cdot \vect {p}$}{k.p} theory}{Doc-Start}{}}
\xrlabel{eq:3.3-kane2@cref}{{[equation][52][3]3.52}{[1][1][]1}}
\xrlabel{eq:3.3-HSO}{{3.53}{}{Spin--orbit coupling}{Doc-Start}{}}
\xrlabel{eq:3.3-HSO@cref}{{[equation][53][3]3.53}{[1][1][]1}}
\xrlabel{eq:3.3-Emin}{{3.57}{}{The minimum pair energy}{Doc-Start}{}}
\xrlabel{eq:3.3-Emin@cref}{{[equation][57][3]3.57}{[1][1][]1}}
\xrlabel{we:3-kc}{{3.5}{}{The minimum pair energy}{Doc-Start}{}}
\xrlabel{we:3-kc@cref}{{[workedexample][5][3]3.5}{[1][1][]1}}
\xrlabel{chap:4}{{4}{}{Linear response}{Doc-Start}{}}
\xrlabel{chap:4@cref}{{[chapter][4][]4}{[1][1][]1}}
\xrlabel{sec:4.1}{{4.1}{}{The Kubo formula}{Doc-Start}{}}
\xrlabel{sec:4.1@cref}{{[section][1][4]4.1}{[1][1][]1}}
\xrlabel{sec:4.2}{{4.2}{}{Correlations and the dynamic structure factor}{Doc-Start}{}}
\xrlabel{sec:4.2@cref}{{[section][2][4]4.2}{[1][1][]1}}
\xrlabel{eq:4.2-Sab}{{4.13}{}{The response matrix}{Doc-Start}{}}
\xrlabel{eq:4.2-Sab@cref}{{[equation][13][4]4.13}{[1][1][]1}}
\xrlabel{eq:4.2-positivity}{{4.20}{}{Positivity}{Doc-Start}{}}
\xrlabel{eq:4.2-positivity@cref}{{[equation][20][4]4.20}{[1][1][]1}}
\xrlabel{sec:4.3}{{4.3}{}{Causality and sum rules}{Doc-Start}{}}
\xrlabel{sec:4.3@cref}{{[section][3][4]4.3}{[1][1][]1}}
\xrlabel{eq:4.3-KK}{{4.24}{}{Kramers--Kronig relations}{Doc-Start}{}}
\xrlabel{eq:4.3-KK@cref}{{[equation][24][4]4.24}{[1][1][]1}}
\xrlabel{eq:4.3-firstmoment}{{4.28}{}{Moments from commutators}{Doc-Start}{}}
\xrlabel{eq:4.3-firstmoment@cref}{{[equation][28][4]4.28}{[1][1][]1}}
\xrlabel{eq:4.3-continuity}{{4.29}{}{The f-sum rule}{Doc-Start}{}}
\xrlabel{eq:4.3-continuity@cref}{{[equation][29][4]4.29}{[1][1][]1}}
\xrlabel{eq:4.3-fsum}{{4.30}{}{The f-sum rule}{Doc-Start}{}}
\xrlabel{eq:4.3-fsum@cref}{{[equation][30][4]4.30}{[1][1][]1}}
\xrlabel{sec:4.4}{{4.4}{}{Screening and the dielectric function}{Doc-Start}{}}
\xrlabel{sec:4.4@cref}{{[section][4][4]4.4}{[1][1][]1}}
\xrlabel{eq:4.4-SL}{{4.45}{}{Structure factor and loss function}{Doc-Start}{}}
\xrlabel{eq:4.4-SL@cref}{{[equation][45][4]4.45}{[1][1][]1}}
\xrlabel{eq:4.4-fsumL}{{4.46}{}{Sum rules for the loss function}{Doc-Start}{}}
\xrlabel{eq:4.4-fsumL@cref}{{[equation][46][4]4.46}{[1][1][]1}}
\xrlabel{eq:4.5-mermin}{{4.62}{}{Collisions and the Mermin function}{Doc-Start}{}}
\xrlabel{eq:4.5-mermin@cref}{{[equation][62][4]4.62}{[1][1][]1}}
\xrlabel{sec:4.6}{{4.6}{}{Local fields and exchange--correlation kernels}{Doc-Start}{}}
\xrlabel{sec:4.6@cref}{{[section][6][4]4.6}{[1][1][]1}}
\xrlabel{eq:4.6-R0}{{4.63}{}{The crystal bubble}{Doc-Start}{}}
\xrlabel{eq:4.6-R0@cref}{{[equation][63][4]4.63}{[1][1][]1}}
\xrlabel{eq:4.6-fxc}{{4.69}{}{Time-dependent density-functional theory}{Doc-Start}{}}
\xrlabel{eq:4.6-fxc@cref}{{[equation][69][4]4.69}{[1][1][]1}}
\xrlabel{eq:4.6-tddft}{{4.70}{}{Time-dependent density-functional theory}{Doc-Start}{}}
\xrlabel{eq:4.6-tddft@cref}{{[equation][70][4]4.70}{[1][1][]1}}
\xrlabel{sec:4.7}{{4.7}{}{Green's functions}{Doc-Start}{}}
\xrlabel{sec:4.7@cref}{{[section][7][4]4.7}{[1][1][]1}}
\xrlabel{sec:4.8}{{4.8}{}{Feynman diagrams}{Doc-Start}{}}
\xrlabel{sec:4.8@cref}{{[section][8][4]4.8}{[1][1][]1}}
\xrlabel{sec:4.9}{{4.9}{}{Excitons and the Bethe--Salpeter equation}{Doc-Start}{}}
\xrlabel{sec:4.9@cref}{{[section][9][4]4.9}{[1][1][]1}}
\xrlabel{chap:5}{{5}{}{Lattice dynamics}{Doc-Start}{}}
\xrlabel{chap:5@cref}{{[chapter][5][]5}{[1][1][]1}}
\xrlabel{sec:5.1}{{5.1}{}{The Born--Oppenheimer approximation}{Doc-Start}{}}
\xrlabel{sec:5.1@cref}{{[section][1][5]5.1}{[1][1][]1}}
\xrlabel{eq:5.1-H}{{5.2}{}{Adiabatic states}{Doc-Start}{}}
\xrlabel{eq:5.1-H@cref}{{[equation][2][5]5.2}{[1][1][]1}}
\xrlabel{eq:5.1-HF}{{5.10}{}{The Hellmann--Feynman force}{Doc-Start}{}}
\xrlabel{eq:5.1-HF@cref}{{[equation][10][5]5.10}{[1][1][]1}}
\xrlabel{fig:5-scales}{{5.2}{}{Scales of electrons and nuclei in silicon. (a) Energies on a logarithmic axis, in eV. Electronic (blue): the atomic unit $\alpha ^2m_e = \SI {27.2}{eV}$, the valence plasma energy $\omega _p = \SI {16.6}{eV}$, the direct gap $E_{\mathrm {dir}} = \SI {3.4}{eV}$ and the indirect gap $\Eg = \SI {1.17}{eV}$. Nuclear (red): the measured zone-center optical phonon, \SI {64.1}{meV} (\SI {15.5}{THz}), the highest vibrational energy of the crystal. Gray dashed: the Born--Oppenheimer estimates $\kappa _{\mathrm {BO}}^n\alpha ^2m_e$ for $n = 2$ (phonons), 3 (cubic anharmonicity), 4 (anharmonic level shifts) and 6 (nonadiabatic level shifts), with $\kappa _{\mathrm {BO}} = (m_e/M)^{1/4} = 0.0665$ for $^{28}$Si. (b) Lengths on a logarithmic axis, in \AA : the lattice constant $a$, the bond length $\sqrt 3a/4$ and the Bohr radius $a_0$ (blue). Red: the root-mean-square displacement of an atom along one Cartesian axis, $\langle u_x^2\rangle ^{1/2}$, at \SI {293}{K} from a lattice-dynamical fit to neutron data and at $T = 0$ from the PBE phonon density of states of the calculation of panel (c) on a $4^3$ grid of phonon wavevectors. Gray dashed: the estimate $\kappa _{\mathrm {BO}}a_0$. (c) Energy per two-atom cell in a PBE calculation with an SG15 norm-conserving pseudopotential, a \SI {50}{Ry} cutoff, an $8^3$ grid of crystal momenta and $a = \SI {5.431}{\angstrom }$. The circles are computed at displacements in which the two atoms move in opposite directions along $[111]$ and the $[111]$ bond changes its length by $u$. The curves are a sixth-degree polynomial fit (solid) and the harmonic term $\tfrac 12\Phi _{\mathrm {opt}}u^2$, $\Phi _{\mathrm {opt}} = \SI {13.6}{eV.\angstrom ^{-2}}$ (dashed), from a quartic fit over $|u| \le \SI {0.1}{\angstrom }$. Shaded: $|u| < \ell _{\mathrm {zp}} = (M_\kappa \omega _{\mathrm {opt}})^{-1/2} = \SI {0.048}{\angstrom }$, the zero-point length of the coordinate $u$. Inset: Hellmann--Feynman force on $u$ (circles) and $-\dd E/\dd u$ from the fit (line)}{Doc-Start}{}}
\xrlabel{fig:5-scales@cref}{{[figure][2][5]5.2}{[1][1][]1}}
\xrlabel{sec:5.2}{{5.2}{}{The harmonic crystal}{Doc-Start}{}}
\xrlabel{sec:5.2@cref}{{[section][2][5]5.2}{[1][1][]1}}
\xrlabel{eq:5.2-expand}{{5.13}{}{Force constants and the dynamical matrix}{Doc-Start}{}}
\xrlabel{eq:5.2-expand@cref}{{[equation][13][5]5.13}{[1][1][]1}}
\xrlabel{eq:5.2-Ddyn}{{5.16}{}{Force constants and the dynamical matrix}{Doc-Start}{}}
\xrlabel{eq:5.2-Ddyn@cref}{{[equation][16][5]5.16}{[1][1][]1}}
\xrlabel{we:5-chain}{{5.2}{}{Force constants and the dynamical matrix}{Doc-Start}{}}
\xrlabel{we:5-chain@cref}{{[workedexample][2][5]5.2}{[1][1][]1}}
\xrlabel{eq:5.2-asr}{{5.18}{}{The acoustic sum rule}{Doc-Start}{}}
\xrlabel{eq:5.2-asr@cref}{{[equation][18][5]5.18}{[1][1][]1}}
\xrlabel{fig:5-si}{{5.3}{}{Phonons of silicon. (a) Dispersion along $[100]$ from $\mathrm {\Gamma }$ to X, along $[110]$ from X$'$ $= (1,1,0)2\pi /a$, which is equivalent to X, through K to $\mathrm {\Gamma }$, and along $[111]$ from $\mathrm {\Gamma }$ to L. Lines: density-functional perturbation theory with the PBE functional, an SG15 norm-conserving pseudopotential, a \SI {60}{Ry} cutoff, a $12^3$ grid of crystal momenta, the measured lattice constant $a = \SI {5.431}{\angstrom }$ and the natural mass of silicon; force constants from a $6^3$ grid of wavevectors, with the acoustic sum rule imposed, interpolated to the path. Circles: frequencies measured by inelastic neutron scattering at room temperature at $\mathrm {\Gamma }$, X and L~\cite {dolling1963inelastic,nilsson1972study}, as tabulated in Ref.~\cite [Table~I]{savrasov1996linear}. Dotted: $\omega = c_\nu Q$ with the sound speeds of the measured elastic constants (\cref {we:5-debye}), drawn up to $Q = 0.25\,(2\pi /a)$. (b) Phonon density of states per two-atom cell, normalized to six modes, from the same force constants on a $48^3$ grid (tetrahedron integration), and the Debye model \cref {eq:5.2-gD} with $\bar c$ from the measured elastic constants, $\omega _{\mathrm {D}} = \SI {55.6}{meV}$, in the convention of one wavevector per atom. One THz corresponds to \SI {4.136}{meV}}{Doc-Start}{}}
\xrlabel{fig:5-si@cref}{{[figure][3][5]5.3}{[1][1][]1}}
\xrlabel{sec:5.3}{{5.3}{}{Phonons}{Doc-Start}{}}
\xrlabel{sec:5.3@cref}{{[section][3][5]5.3}{[1][1][]1}}
\xrlabel{sec:5.5}{{5.4}{}{Polar crystals}{Doc-Start}{}}
\xrlabel{sec:5.5@cref}{{[section][4][5]5.4}{[1][1][]1}}
\xrlabel{sec:5.6}{{5.5}{}{Displacement correlations}{Doc-Start}{}}
\xrlabel{sec:5.6@cref}{{[section][5][5]5.5}{[1][1][]1}}
\xrlabel{sec:5.7}{{5.6}{}{Anharmonicity}{Doc-Start}{}}
\xrlabel{sec:5.7@cref}{{[section][6][5]5.6}{[1][1][]1}}
\xrlabel{sec:6.2}{{6.2}{}{Electron--phonon matrix elements}{Doc-Start}{}}
\xrlabel{sec:6.2@cref}{{[section][2][6]6.2}{[1][1][]1}}
\xrlabel{eq:6.2-asr}{{6.12}{}{Acoustic phonons and deformation potentials}{Doc-Start}{}}
\xrlabel{eq:6.2-asr@cref}{{[equation][12][6]6.12}{[1][1][]1}}
\xrlabel{eq:6.2-D}{{6.13}{}{Acoustic phonons and deformation potentials}{Doc-Start}{}}
\xrlabel{eq:6.2-D@cref}{{[equation][13][6]6.13}{[1][1][]1}}
\xrlabel{sec:6.3}{{6.3}{}{Electron self-energy from phonons}{Doc-Start}{}}
\xrlabel{sec:6.3@cref}{{[section][3][6]6.3}{[1][1][]1}}
\xrlabel{sec:6.4}{{6.4}{}{Phonon self-energy and the Eliashberg function}{Doc-Start}{}}
\xrlabel{sec:6.4@cref}{{[section][4][6]6.4}{[1][1][]1}}
\xrlabel{sec:6.5}{{6.5}{}{Phonon-assisted optical absorption}{Doc-Start}{}}
\xrlabel{sec:6.5@cref}{{[section][5][6]6.5}{[1][1][]1}}
\xrlabel{sec:7.1}{{7.1}{}{Magnetic order and spin waves}{Doc-Start}{}}
\xrlabel{sec:7.1@cref}{{[section][1][7]7.1}{[1][1][]1}}
\xrlabel{eq:7.1-heis}{{7.6}{}{The Heisenberg model}{Doc-Start}{}}
\xrlabel{eq:7.1-heis@cref}{{[equation][6][7]7.6}{[1][1][]1}}
\xrlabel{sec:7.2}{{7.2}{}{Superconductivity}{Doc-Start}{}}
\xrlabel{sec:7.2@cref}{{[section][2][7]7.2}{[1][1][]1}}
\xrlabel{sec:7.3}{{7.3}{}{Dirac and narrow-gap materials}{Doc-Start}{}}
\xrlabel{sec:7.3@cref}{{[section][3][7]7.3}{[1][1][]1}}
\xrlabel{eq:7.3-alpha}{{7.46}{}{Screening by Dirac electrons}{Doc-Start}{}}
\xrlabel{eq:7.3-alpha@cref}{{[equation][46][7]7.46}{[1][1][]1}}
\xrlabel{sec:7.4}{{7.4}{}{Defects and dopants}{Doc-Start}{}}
\xrlabel{sec:7.4@cref}{{[section][4][7]7.4}{[1][1][]1}}
\xrlabel{chap:8}{{8}{}{The response framework}{Doc-Start}{}}
\xrlabel{chap:8@cref}{{[chapter][8][]8}{[1][1][]1}}
\xrlabel{sec:8.1}{{8.1}{}{Dark-matter couplings to a target}{Doc-Start}{}}
\xrlabel{sec:8.1@cref}{{[section][1][8]8.1}{[1][1][]1}}
\xrlabel{eq:8.1-rows}{{8.3}{}{Target operators per species}{Doc-Start}{}}
\xrlabel{eq:8.1-rows@cref}{{[equation][3][8]8.3}{[1][1][]1}}
\xrlabel{tab:8-operators}{{8.1}{}{Non-relativistic operators of dark-matter scattering and absorption, and the target operators each needs. Upper block: the operators $\mathcal {O}_i$ of Refs.~\cite [Eqs.~(19)--(24) and (45)]{fitzpatrick2013effective} and~\cite [Eqs.~(12) and (13)]{anand2014model}, written for an electron with $\vect {q}$ the momentum the dark matter loses; $\vect {v}^\perp $ is \cref {eq:8.1-vperp}, $\bar {\vect {v}} \equiv \bar {\vect {v}}_\chi = \vect {v} - \vect {q}/2\mchi $, and $m = m_e$. The rows are those of \cref {eq:8.1-rows} for $s = e$, all at the same $\vect {q}$; repeated Cartesian indices are summed, and $\epsilon _{\alpha \beta \gamma }$ is the Levi-Civita symbol. For nucleons, replace $e$ by $N$ and $m_e$ by the nucleon mass; the nuclear rows follow from \cref {eq:8.1-nuclear}. Lower block: absorption of a boson $\phi $ of mass $m_\phi $ at $\vect {q} = \vect {p}_\phi $, and fermionic absorption with an outgoing fermion $\tilde \nu $ (\cref {sec:8.1-absorption}). Last column: what each row excites in an insulator. Through the nuclei the density row also excites phonons, and in an ordered magnet the spin rows excite magnons (\cref {sec:7.1})}{Doc-Start}{}}
\xrlabel{tab:8-operators@cref}{{[table][1][8]8.1}{[1][1][]1}}
\xrlabel{eq:8.1-nuclear}{{8.8}{}{Nuclei}{Doc-Start}{}}
\xrlabel{eq:8.1-nuclear@cref}{{[equation][8][8]8.8}{[1][1][]1}}
\xrlabel{sec:8.2}{{8.2}{}{The master rate formula}{Doc-Start}{}}
\xrlabel{sec:8.2@cref}{{[section][2][8]8.2}{[1][1][]1}}
\xrlabel{sec:8.3}{{8.3}{}{Exact relations}{Doc-Start}{}}
\xrlabel{sec:8.3@cref}{{[section][3][8]8.3}{[1][1][]1}}
\xrlabel{eq:8.3-hankel}{{8.23}{}{Moment matrices}{Doc-Start}{}}
\xrlabel{eq:8.3-hankel@cref}{{[equation][23][8]8.23}{[1][1][]1}}
\xrlabel{eq:8.3-ward}{{8.29}{}{Translation Ward identities}{Doc-Start}{}}
\xrlabel{eq:8.3-ward@cref}{{[equation][29][8]8.29}{[1][1][]1}}
\xrlabel{eq:8.3-soft}{{8.31}{}{Translation Ward identities}{Doc-Start}{}}
\xrlabel{eq:8.3-soft@cref}{{[equation][31][8]8.31}{[1][1][]1}}
\xrlabel{eq:8.3-boostward}{{8.32}{}{Translation Ward identities}{Doc-Start}{}}
\xrlabel{eq:8.3-boostward@cref}{{[equation][32][8]8.32}{[1][1][]1}}
\xrlabel{sec:8.4}{{8.4}{}{Screening for every coupling}{Doc-Start}{}}
\xrlabel{sec:8.4@cref}{{[section][4][8]8.4}{[1][1][]1}}
\xrlabel{sec:8.5}{{8.5}{}{The expansion hierarchy}{Doc-Start}{}}
\xrlabel{sec:8.5@cref}{{[section][5][8]8.5}{[1][1][]1}}
\xrlabel{eq:8.5-label}{{8.44}{}{The truncation label}{Doc-Start}{}}
\xrlabel{eq:8.5-label@cref}{{[equation][44][8]8.44}{[1][1][]1}}
\xrlabel{sec:8.6}{{8.6}{}{The decomposition theorem}{Doc-Start}{}}
\xrlabel{sec:8.6@cref}{{[section][6][8]8.6}{[1][1][]1}}
\xrlabel{eq:8.6-fsum-e}{{8.54}{}{Partial sum rules}{Doc-Start}{}}
\xrlabel{eq:8.6-fsum-e@cref}{{[equation][54][8]8.54}{[1][1][]1}}
\xrlabel{eq:8.6-fsum-k}{{8.56}{}{Partial sum rules}{Doc-Start}{}}
\xrlabel{eq:8.6-fsum-k@cref}{{[equation][56][8]8.56}{[1][1][]1}}
\xrlabel{eq:8.6-fsum-ek}{{8.57}{}{Partial sum rules}{Doc-Start}{}}
\xrlabel{eq:8.6-fsum-ek@cref}{{[equation][57][8]8.57}{[1][1][]1}}
\xrlabel{sec:8.7}{{8.7}{}{Final-state-resolved response}{Doc-Start}{}}
\xrlabel{sec:8.7@cref}{{[section][7][8]8.7}{[1][1][]1}}
\xrlabel{eq:8.7-screened}{{8.68}{}{Pairs in an insulator}{Doc-Start}{}}
\xrlabel{eq:8.7-screened@cref}{{[equation][68][8]8.68}{[1][1][]1}}
\xrlabel{we:8-pairs}{{8.9}{}{Pairs in an insulator}{Doc-Start}{}}
\xrlabel{we:8-pairs@cref}{{[workedexample][9][8]8.9}{[1][1][]1}}
\xrlabel{sec:8.8}{{8.8}{}{Uncertainty structure}{Doc-Start}{}}
\xrlabel{sec:8.8@cref}{{[section][8][8]8.8}{[1][1][]1}}
\xrlabel{eq:8.8-delta}{{8.71}{}{The nuisance of a truncation}{Doc-Start}{}}
\xrlabel{eq:8.8-delta@cref}{{[equation][71][8]8.71}{[1][1][]1}}
\xrlabel{eq:8.8-geodesic}{{8.73}{}{The nuisance of a truncation}{Doc-Start}{}}
\xrlabel{eq:8.8-geodesic@cref}{{[equation][73][8]8.73}{[1][1][]1}}
\xrlabel{eq:8.8-open}{{8.75}{}{Scale, tilt and edge shifts}{Doc-Start}{}}
\xrlabel{eq:8.8-open@cref}{{[equation][75][8]8.75}{[1][1][]1}}
\xrlabel{eq:8.8-fsum}{{8.76}{}{Exact relations as constraints}{Doc-Start}{}}
\xrlabel{eq:8.8-fsum@cref}{{[equation][76][8]8.76}{[1][1][]1}}
\xrlabel{app:C}{{C}{}{Mathematical tools}{Doc-Start}{}}
\xrlabel{app:C@cref}{{[appendix][3][2147483647]C}{[1][1][]1}}
\xrlabel{app:C.2}{{C.2}{}{Box normalization and Brillouin-zone sums}{Doc-Start}{}}
\xrlabel{app:C.2@cref}{{[subappendix][2][2147483647,3]C.2}{[1][1][]1}}
\xrlabel{eq:C.2-grid}{{C.15}{}{Zone averages and grids}{Doc-Start}{}}
\xrlabel{eq:C.2-grid@cref}{{[equation][15][2147483647,3]C.15}{[1][1][]1}}
\xrlabel{app:C.3}{{C.3}{}{Angular averages and spherical functions}{Doc-Start}{}}
\xrlabel{app:C.3@cref}{{[subappendix][3][2147483647,3]C.3}{[1][1][]1}}
\xrlabel{eq:C.3-planewave}{{C.22}{}{The plane-wave expansion}{Doc-Start}{}}
\xrlabel{eq:C.3-planewave@cref}{{[equation][22][2147483647,3]C.22}{[1][1][]1}}
\xrlabel{app:C.5}{{C.5}{}{Complex analysis for response functions}{Doc-Start}{}}
\xrlabel{app:C.5@cref}{{[subappendix][5][2147483647,3]C.5}{[1][1][]1}}
\xrlabel{app:C.8}{{C.8}{}{Special functions}{Doc-Start}{}}
\xrlabel{app:C.8@cref}{{[subappendix][8][2147483647,3]C.8}{[1][1][]1}}
\xrlabel{eq:C.8-FL}{{C.72}{}{The Lindhard logarithms}{Doc-Start}{}}
\xrlabel{eq:C.8-FL@cref}{{[equation][72][2147483647,3]C.72}{[1][1][]1}}
\xrlabel{app:C.9}{{C.9}{}{Classical fields and static propagators}{Doc-Start}{}}
\xrlabel{app:C.9@cref}{{[subappendix][9][2147483647,3]C.9}{[1][1][]1}}
\xrlabel{eq:C.9-Gq}{{C.79}{}{The static Green's function}{Doc-Start}{}}
\xrlabel{eq:C.9-Gq@cref}{{[equation][79][2147483647,3]C.79}{[1][1][]1}}
\xrlabel{eq:C.9-mix}{{C.83}{}{Kinetic mixing}{Doc-Start}{}}
\xrlabel{eq:C.9-mix@cref}{{[equation][83][2147483647,3]C.83}{[1][1][]1}}
\xrlabel{eq:C.12-sip}{{C.109}{}{Semi-infinite programs}{Doc-Start}{}}
\xrlabel{eq:C.12-sip@cref}{{[equation][109][2147483647,3]C.109}{[1][1][]1}}
\xrlabel{app:C.15}{{C.15}{}{Dirac equation and Foldy--Wouthuysen reduction}{Doc-Start}{}}
\xrlabel{app:C.15@cref}{{[subappendix][15][2147483647,3]C.15}{[1][1][]1}}
\xrlabel{app:C.17}{{C.17}{}{Point and space groups}{Doc-Start}{}}
\xrlabel{app:C.17@cref}{{[subappendix][17][2147483647,3]C.17}{[1][1][]1}}
\xrlabel{app:C.18}{{C.18}{}{Continuum elasticity}{Doc-Start}{}}
\xrlabel{app:C.18@cref}{{[subappendix][18][2147483647,3]C.18}{[1][1][]1}}
\xrlabel{app:C.20}{{C.20}{}{Dissipative two-level systems and kinks}{Doc-Start}{}}
\xrlabel{app:C.20@cref}{{[subappendix][20][2147483647,3]C.20}{[1][1][]1}}
\xrlabel{app:D}{{D}{}{Long derivations}{Doc-Start}{}}
\xrlabel{app:D@cref}{{[appendix][4][2147483647]D}{[1][1][]1}}
\xrlabel{app:D.6}{{D.4}{}{The f-sum rule}{Doc-Start}{}}
\xrlabel{app:D.6@cref}{{[subappendix][4][2147483647,4]D.4}{[1][1][]1}}
\xrlabel{app:D.11}{{D.9}{}{Wick's theorem and the Feynman rules}{Doc-Start}{}}
\xrlabel{app:D.11@cref}{{[subappendix][9][2147483647,4]D.9}{[1][1][]1}}
\xrlabel{app:D.23}{{D.17}{}{Effective-mass theory of shallow donors}{Doc-Start}{}}
\xrlabel{app:D.23@cref}{{[subappendix][17][2147483647,4]D.17}{[1][1][]1}}
\xrlabel{chap:9}{{9}{}{Channel atlas}{Doc-Start}{}}
\xrlabel{chap:9@cref}{{[chapter][9][]9}{[1][1][]1}}
\xrlabel{sec:9.2}{{9.2}{}{The kinematic map}{Doc-Start}{}}
\xrlabel{sec:9.2@cref}{{[section][2][9]9.2}{[1][1][]1}}
\xrlabel{sec:9.4}{{9.4}{}{Limit Checks and the truncation label}{Doc-Start}{}}
\xrlabel{sec:9.4@cref}{{[section][4][9]9.4}{[1][1][]1}}
\xrlabel{chap:10}{{10}{}{Electronic excitations}{Doc-Start}{}}
\xrlabel{chap:10@cref}{{[chapter][10][]10}{[1][1][]1}}
\xrlabel{sec:10.1}{{10.1}{}{The crystal form factor}{Doc-Start}{}}
\xrlabel{sec:10.1@cref}{{[section][1][10]10.1}{[1][1][]1}}
\xrlabel{sec:10.2}{{10.2}{}{Screening and the loss-function rate}{Doc-Start}{}}
\xrlabel{sec:10.2@cref}{{[section][2][10]10.2}{[1][1][]1}}
\xrlabel{sec:10.4}{{10.4}{}{Spin and current couplings}{Doc-Start}{}}
\xrlabel{sec:10.4@cref}{{[section][4][10]10.4}{[1][1][]1}}
\xrlabel{tab:10-minimal}{{10.1}{}{The minimal response table: the rows of the electronic response matrix that each non-relativistic operator of \cref {tab:8-operators} needs, and their screening. The longitudinal current and spin current have been eliminated by continuity (elastic scattering, local potentials, electron spin conserved). Screening is for a cubic crystal without spin--orbit coupling or magnetic order. Singlet rows make spin-singlet electron--hole pairs and triplet rows spin-triplet pairs; the two blocks do not interfere. Notation: $j_{\mathrm {T}}$, current across $\vect {q}$; $\mathcal {J}_{\mathrm {T}}$: spin current with its velocity index across $\vect {q}$. With inversion and time reversal, the singlet block on $\{n_e, j_{\mathrm {T}1}, j_{\mathrm {T}2}\}$ holds six real functions per $(\vect {q},\omega )$. For independent particles the triplet block is one quarter of it (\cref {eq:D.30-quarter})}{Doc-Start}{}}
\xrlabel{tab:10-minimal@cref}{{[table][1][10]10.1}{[1][1][]1}}
\xrlabel{eq:10.4-band}{{10.15}{}{Interference bound on a mixed block}{Doc-Start}{}}
\xrlabel{eq:10.4-band@cref}{{[equation][15][10]10.15}{[1][1][]1}}
\xrlabel{sec:10.5}{{10.5}{}{Excitons at finite momentum}{Doc-Start}{}}
\xrlabel{sec:10.5@cref}{{[section][5][10]10.5}{[1][1][]1}}
\xrlabel{eq:10.5-res}{{10.22}{}{Time orderings and the coupling block}{Doc-Start}{}}
\xrlabel{eq:10.5-res@cref}{{[equation][22][10]10.22}{[1][1][]1}}
\xrlabel{sec:10.6}{{10.6}{}{Core and relativistic states}{Doc-Start}{}}
\xrlabel{sec:10.6@cref}{{[section][6][10]10.6}{[1][1][]1}}
\xrlabel{sec:10.7}{{10.7}{}{Absorption on electrons}{Doc-Start}{}}
\xrlabel{sec:10.7@cref}{{[section][7][10]10.7}{[1][1][]1}}
\xrlabel{eq:10.7-alp}{{10.33}{}{Axion-like-particle absorption}{Doc-Start}{}}
\xrlabel{eq:10.7-alp@cref}{{[equation][33][10]10.33}{[1][1][]1}}
\xrlabel{sec:10.7-scalar}{{10.7.4}{}{Scalar absorption and the kinetic-energy row}{Doc-Start}{}}
\xrlabel{sec:10.7-scalar@cref}{{[subsection][4][10,7]10.7.4}{[1][1][]1}}
\xrlabel{eq:10.7-scalar}{{10.36}{}{Scalar absorption and the kinetic-energy row}{Doc-Start}{}}
\xrlabel{eq:10.7-scalar@cref}{{[equation][36][10]10.36}{[1][1][]1}}
\xrlabel{sec:10.8}{{10.8}{}{Two-insertion processes}{Doc-Start}{}}
\xrlabel{sec:10.8@cref}{{[section][8][10]10.8}{[1][1][]1}}
\xrlabel{eq:10.8-P}{{10.44}{}{Pair absorption}{Doc-Start}{}}
\xrlabel{eq:10.8-P@cref}{{[equation][44][10]10.44}{[1][1][]1}}
\xrlabel{sec:10.9}{{10.9}{}{Inelastic and relativistic dark matter}{Doc-Start}{}}
\xrlabel{sec:10.9@cref}{{[section][9][10]10.9}{[1][1][]1}}
\xrlabel{sec:10.10}{{10.10}{}{Nuisance parameters of the electronic rate}{Doc-Start}{}}
\xrlabel{sec:10.10@cref}{{[section][10][10]10.10}{[1][1][]1}}
\xrlabel{chap:11}{{11}{}{Lattice excitations}{Doc-Start}{}}
\xrlabel{chap:11@cref}{{[chapter][11][]11}{[1][1][]1}}
\xrlabel{sec:11.1}{{11.1}{}{Acoustic phonons}{Doc-Start}{}}
\xrlabel{sec:11.1@cref}{{[section][1][11]11.1}{[1][1][]1}}
\xrlabel{sec:11.2}{{11.2}{}{Optical phonons in polar crystals}{Doc-Start}{}}
\xrlabel{sec:11.2@cref}{{[section][2][11]11.2}{[1][1][]1}}
\xrlabel{sec:11.3}{{11.3}{}{Multiphonon response}{Doc-Start}{}}
\xrlabel{sec:11.3@cref}{{[section][3][11]11.3}{[1][1][]1}}
\xrlabel{eq:11.3-Bkk}{{11.24}{}{Exact phonon-order weights}{Doc-Start}{}}
\xrlabel{eq:11.3-Bkk@cref}{{[equation][24][11]11.24}{[1][1][]1}}
\xrlabel{eq:11.3-weight}{{11.25}{}{Exact phonon-order weights}{Doc-Start}{}}
\xrlabel{eq:11.3-weight@cref}{{[equation][25][11]11.25}{[1][1][]1}}
\xrlabel{sec:11.4}{{11.4}{}{Impulse limit and free nuclear recoil}{Doc-Start}{}}
\xrlabel{sec:11.4@cref}{{[section][4][11]11.4}{[1][1][]1}}
\xrlabel{sec:11.5}{{11.5}{}{Anharmonic and quantum-crystal corrections}{Doc-Start}{}}
\xrlabel{sec:11.5@cref}{{[section][5][11]11.5}{[1][1][]1}}
\xrlabel{sec:11.6}{{11.6}{}{Nuclear-spin and isotope couplings}{Doc-Start}{}}
\xrlabel{sec:11.6@cref}{{[section][6][11]11.6}{[1][1][]1}}
\xrlabel{eq:11.6-spinS}{{11.38}{}{Unpolarized nuclear spins}{Doc-Start}{}}
\xrlabel{eq:11.6-spinS@cref}{{[equation][38][11]11.38}{[1][1][]1}}
\xrlabel{eq:11.6-lifshitz}{{11.40}{}{An isotopic mass defect}{Doc-Start}{}}
\xrlabel{eq:11.6-lifshitz@cref}{{[equation][40][11]11.40}{[1][1][]1}}
\xrlabel{eq:11.6-Tbound}{{11.44}{}{Absorption of a baryon-minus-lepton vector}{Doc-Start}{}}
\xrlabel{eq:11.6-Tbound@cref}{{[equation][44][11]11.44}{[1][1][]1}}
\xrlabel{sec:11.7}{{11.7}{}{Anisotropy and daily modulation}{Doc-Start}{}}
\xrlabel{sec:11.7@cref}{{[section][7][11]11.7}{[1][1][]1}}
\xrlabel{sec:11.8}{{11.8}{}{Absorption on phonons and acoustic modes}{Doc-Start}{}}
\xrlabel{sec:11.8@cref}{{[section][8][11]11.8}{[1][1][]1}}
\xrlabel{eq:11.8-h}{{11.59}{}{Resonant-mass detectors}{Doc-Start}{}}
\xrlabel{eq:11.8-h@cref}{{[equation][59][11]11.59}{[1][1][]1}}
\xrlabel{sec:11.9}{{11.9}{}{Nuisance parameters of the lattice rates}{Doc-Start}{}}
\xrlabel{sec:11.9@cref}{{[section][9][11]11.9}{[1][1][]1}}
\xrlabel{sec:12.1}{{12.1}{}{The joint electron--lattice response}{Doc-Start}{}}
\xrlabel{sec:12.1@cref}{{[section][1][12]12.1}{[1][1][]1}}
\xrlabel{eq:12.1-Skk}{{12.6}{}{The boosted electronic response}{Doc-Start}{}}
\xrlabel{eq:12.1-Skk@cref}{{[equation][6][12]12.6}{[1][1][]1}}
\xrlabel{tab:12-truncations}{{12.1}{}{The published Migdal formulas as truncations of the Migdal entry \cref {eq:12.1-Skk} of the joint response. Each keeps the boost $\vect {q}\cdot \vect {p}/(M_\kappa \omega _e^2)$ at Born--Oppenheimer order $n_B = 2$; the columns give the choices for the vertex and electronic factor, for the lattice factor, and the section that treats the form}{Doc-Start}{}}
\xrlabel{tab:12-truncations@cref}{{[table][1][12]12.1}{[1][1][]1}}
\xrlabel{sec:12.2}{{12.2}{}{The atomic limit}{Doc-Start}{}}
\xrlabel{sec:12.2@cref}{{[section][2][12]12.2}{[1][1][]1}}
\xrlabel{sec:12.3}{{12.3}{}{The loss-function Migdal probability}{Doc-Start}{}}
\xrlabel{sec:12.3@cref}{{[section][3][12]12.3}{[1][1][]1}}
\xrlabel{sec:12.3-validity}{{12.3.2}{}{Validity of the soft limit}{Doc-Start}{}}
\xrlabel{sec:12.3-validity@cref}{{[subsection][2][12,3]12.3.2}{[1][1][]1}}
\xrlabel{lim:12.3a}{{12.2}{}{Validity of the soft limit}{Doc-Start}{}}
\xrlabel{lim:12.3a@cref}{{[tcb@cnt@limitbox][2][12]12.2}{[1][1][]1}}
\xrlabel{sec:12.4}{{12.4}{}{Phonon final states at small momentum}{Doc-Start}{}}
\xrlabel{sec:12.4@cref}{{[section][4][12]12.4}{[1][1][]1}}
\xrlabel{sec:12.5}{{12.5}{}{Moments of the Migdal spectrum}{Doc-Start}{}}
\xrlabel{sec:12.5@cref}{{[section][5][12]12.5}{[1][1][]1}}
\xrlabel{eq:12.5-trace}{{12.39}{}{On-site drag mass and trace rule}{Doc-Start}{}}
\xrlabel{eq:12.5-trace@cref}{{[equation][39][12]12.39}{[1][1][]1}}
\xrlabel{sec:12.5-bounds}{{12.5.3}{}{Bounds on the Migdal probability}{Doc-Start}{}}
\xrlabel{sec:12.5-bounds@cref}{{[subsection][3][12,5]12.5.3}{[1][1][]1}}
\xrlabel{eq:12.5-Rrel}{{12.44}{}{Third moment and contact density}{Doc-Start}{}}
\xrlabel{eq:12.5-Rrel@cref}{{[equation][44][12]12.44}{[1][1][]1}}
\xrlabel{sec:12.6}{{12.6}{}{Bremsstrahlung and Migdal-type absorption}{Doc-Start}{}}
\xrlabel{sec:12.6@cref}{{[section][6][12]12.6}{[1][1][]1}}
\xrlabel{eq:12.6-accel}{{12.50}{}{Migdal-type absorption of a scalar}{Doc-Start}{}}
\xrlabel{eq:12.6-accel@cref}{{[equation][50][12]12.50}{[1][1][]1}}
\xrlabel{eq:12.6-absrate}{{12.52}{}{Migdal-type absorption of a scalar}{Doc-Start}{}}
\xrlabel{eq:12.6-absrate@cref}{{[equation][52][12]12.52}{[1][1][]1}}
\xrlabel{tab:12-nuisances}{{12.3}{}{Nuisances of the Migdal sector. Type in the classes of \cref {sec:8.8}: normalization, tilt (in $\omega _e$, in $q$ or in the dark-matter mass), threshold and lattice spectrum (the share of a Migdal event that goes to phonons, at fixed electronic energy transfer). Basis of the size: ``computed'' (next order computed), ``model'' (next order computed in a model), ``est.''\ (parametric estimate), ``bound'' (only an inequality is known), ``named'' (no size). Sizes are for silicon unless another crystal is named. The last column gives the section where the item is treated}{Doc-Start}{}}
\xrlabel{tab:12-nuisances@cref}{{[table][3][12]12.3}{[1][1][]1}}
\xrlabel{sec:12.6-open}{{12.6.5}{}{Open problems}{Doc-Start}{}}
\xrlabel{sec:12.6-open@cref}{{[subsection][5][12,6]12.6.5}{[1][1][]1}}
\xrlabel{chap:13}{{13}{}{Phonon-assisted electron recoil}{Doc-Start}{}}
\xrlabel{chap:13@cref}{{[chapter][13][]13}{[1][1][]1}}
\xrlabel{sec:13.1}{{13.1}{}{The indirect-gap window}{Doc-Start}{}}
\xrlabel{sec:13.1@cref}{{[section][1][13]13.1}{[1][1][]1}}
\xrlabel{sec:13.2}{{13.2}{}{The phonon-assisted amplitude}{Doc-Start}{}}
\xrlabel{sec:13.2@cref}{{[section][2][13]13.2}{[1][1][]1}}
\xrlabel{sec:13.3}{{13.3}{}{The phonon-assisted rate}{Doc-Start}{}}
\xrlabel{sec:13.3@cref}{{[section][3][13]13.3}{[1][1][]1}}
\xrlabel{sec:13.4}{{13.4}{}{Screening of the phonon-assisted rate}{Doc-Start}{}}
\xrlabel{sec:13.4@cref}{{[section][4][13]13.4}{[1][1][]1}}
\xrlabel{sec:13.5}{{13.5}{}{The optical limit}{Doc-Start}{}}
\xrlabel{sec:13.5@cref}{{[section][5][13]13.5}{[1][1][]1}}
\xrlabel{sec:13.6}{{13.6}{}{The off-shell window}{Doc-Start}{}}
\xrlabel{sec:13.6@cref}{{[section][6][13]13.6}{[1][1][]1}}
\xrlabel{sec:13.7}{{13.7}{}{Temperature and zero-point renormalization}{Doc-Start}{}}
\xrlabel{sec:13.7@cref}{{[section][7][13]13.7}{[1][1][]1}}
\xrlabel{eq:13.7-fsum}{{13.36}{}{The renormalized zero-phonon response}{Doc-Start}{}}
\xrlabel{eq:13.7-fsum@cref}{{[equation][36][13]13.36}{[1][1][]1}}
\xrlabel{sec:13.8}{{13.8}{}{Exciton replicas and bosonic absorption}{Doc-Start}{}}
\xrlabel{sec:13.8@cref}{{[section][8][13]13.8}{[1][1][]1}}
\xrlabel{eq:13.8-MT}{{13.48}{}{Phonon-assisted absorption of bosonic dark matter}{Doc-Start}{}}
\xrlabel{eq:13.8-MT@cref}{{[equation][48][13]13.48}{[1][1][]1}}
\xrlabel{eq:13.8-scalar}{{13.49}{}{Phonon-assisted absorption of bosonic dark matter}{Doc-Start}{}}
\xrlabel{eq:13.8-scalar@cref}{{[equation][49][13]13.49}{[1][1][]1}}
\xrlabel{sec:13.9}{{13.9}{}{The silicon window}{Doc-Start}{}}
\xrlabel{sec:13.9@cref}{{[section][9][13]13.9}{[1][1][]1}}
\xrlabel{eq:13.9-glide}{{13.50}{}{The glide-parity selection rule}{Doc-Start}{}}
\xrlabel{eq:13.9-glide@cref}{{[equation][50][13]13.50}{[1][1][]1}}
\xrlabel{sec:13.9-rates}{{13.9.4}{}{Rates across the window}{Doc-Start}{}}
\xrlabel{sec:13.9-rates@cref}{{[subsection][4][13,9]13.9.4}{[1][1][]1}}
\xrlabel{sec:13.9-open}{{13.9.8}{}{Open problems}{Doc-Start}{}}
\xrlabel{sec:13.9-open@cref}{{[subsection][8][13,9]13.9.8}{[1][1][]1}}
\xrlabel{chap:14}{{14}{}{Magnetic excitations}{Doc-Start}{}}
\xrlabel{chap:14@cref}{{[chapter][14][]14}{[1][1][]1}}
\xrlabel{sec:14.1}{{14.1}{}{Magnon scattering}{Doc-Start}{}}
\xrlabel{sec:14.1@cref}{{[section][1][14]14.1}{[1][1][]1}}
\xrlabel{sec:14.2}{{14.2}{}{Two-magnon continua and magnon--phonon modes}{Doc-Start}{}}
\xrlabel{sec:14.2@cref}{{[section][2][14]14.2}{[1][1][]1}}
\xrlabel{eq:14.2-ladder}{{14.12}{}{Final-state interactions and bound pairs}{Doc-Start}{}}
\xrlabel{eq:14.2-ladder@cref}{{[equation][12][14]14.12}{[1][1][]1}}
\xrlabel{sec:14.3}{{14.3}{}{Absorption into magnons}{Doc-Start}{}}
\xrlabel{sec:14.3@cref}{{[section][3][14]14.3}{[1][1][]1}}
\xrlabel{sec:14.5}{{14.5}{}{Chiral and nonreciprocal targets}{Doc-Start}{}}
\xrlabel{sec:14.5@cref}{{[section][5][14]14.5}{[1][1][]1}}
\xrlabel{eq:14.5-GOA}{{14.35}{}{Reciprocity of the response matrix}{Doc-Start}{}}
\xrlabel{eq:14.5-GOA@cref}{{[equation][35][14]14.35}{[1][1][]1}}
\xrlabel{sec:14.6}{{14.6}{}{Coherent fields and precessing nuclear spins}{Doc-Start}{}}
\xrlabel{sec:14.6@cref}{{[section][6][14]14.6}{[1][1][]1}}
\xrlabel{chap:15}{{15}{}{Superconductors}{Doc-Start}{}}
\xrlabel{chap:15@cref}{{[chapter][15][]15}{[1][1][]1}}
\xrlabel{sec:15.1}{{15.1}{}{Pair breaking by dark matter}{Doc-Start}{}}
\xrlabel{sec:15.1@cref}{{[section][1][15]15.1}{[1][1][]1}}
\xrlabel{sec:15.2}{{15.2}{}{Screening at the kinematic edge}{Doc-Start}{}}
\xrlabel{sec:15.2@cref}{{[section][2][15]15.2}{[1][1][]1}}
\xrlabel{tab:15-edge}{{15.1}{}{Loss function of aluminum on the kinematic edge $q = \omega /v_{\mathrm {max}}$ ($\omega /qv_{\mathrm {F}} = v_{\mathrm {max}}/v_{\mathrm {F}} = 0.392$), divided by $\mathfrak {g}_{\mathrm {I}}L^{\mathrm {nor}}$, the value with normal-state screening, against $y = \omega /\Delta _{\mathrm {BCS}}$. Columns: case-I ratio $\mathfrak {g}_{\mathrm {I}}$ of \cref {eq:15.1-gI}; vertex fraction $\mathfrak {v}$ of \cref {sec:15.1-bare}; conserving response \cref {eq:15.1-Rcons}, clean, and with Born scattering by impurities at the mean free paths $\mathfrak {l}_e = \SI {1}{\micro m}$ and \SI {30}{nm} (\cref {sec:15.3-sc}), each divided by $\mathfrak {g}_{\mathrm {I}}$ times the normal state at the same scattering rate. All at $T = 0$, constant density of states, $q \ll k_{\mathrm {F}}$}{Doc-Start}{}}
\xrlabel{tab:15-edge@cref}{{[table][1][15]15.1}{[1][1][]1}}
\xrlabel{sec:15.3}{{15.3}{}{Films and disorder}{Doc-Start}{}}
\xrlabel{sec:15.3@cref}{{[section][3][15]15.3}{[1][1][]1}}
\xrlabel{sec:15.4}{{15.4}{}{Absorption in superconductors and metals}{Doc-Start}{}}
\xrlabel{sec:15.4@cref}{{[section][4][15]15.4}{[1][1][]1}}
\xrlabel{eq:15.4-holstein}{{15.24}{}{Phonon-assisted absorption}{Doc-Start}{}}
\xrlabel{eq:15.4-holstein@cref}{{[equation][24][15]15.24}{[1][1][]1}}
\xrlabel{sec:15.4-spin}{{15.4.3}{}{Spin-current absorption}{Doc-Start}{}}
\xrlabel{sec:15.4-spin@cref}{{[subsection][3][15,4]15.4.3}{[1][1][]1}}
\xrlabel{sec:15.4-strong}{{15.4.4}{}{Strong coupling and the transport function}{Doc-Start}{}}
\xrlabel{sec:15.4-strong@cref}{{[subsection][4][15,4]15.4.4}{[1][1][]1}}
\xrlabel{chap:16}{{16}{}{Novel targets}{Doc-Start}{}}
\xrlabel{chap:16@cref}{{[chapter][16][]16}{[1][1][]1}}
\xrlabel{sec:16.1}{{16.1}{}{Dirac-material rates}{Doc-Start}{}}
\xrlabel{sec:16.1@cref}{{[section][1][16]16.1}{[1][1][]1}}
\xrlabel{sec:16.1-zrte5}{{16.1.3}{}{ZrTe\texorpdfstring {$_5$}{5} and other anisotropic targets}{Doc-Start}{}}
\xrlabel{sec:16.1-zrte5@cref}{{[subsection][3][16,1]16.1.3}{[1][1][]1}}
\xrlabel{sec:16.2}{{16.2}{}{Low-dimensional plasmons and layered stacks}{Doc-Start}{}}
\xrlabel{sec:16.2@cref}{{[section][2][16]16.2}{[1][1][]1}}
\xrlabel{sec:16.3}{{16.3}{}{Graphene, nanotubes and quantum dots}{Doc-Start}{}}
\xrlabel{sec:16.3@cref}{{[section][3][16]16.3}{[1][1][]1}}
\xrlabel{sec:16.4}{{16.4}{}{Molecular crystals}{Doc-Start}{}}
\xrlabel{sec:16.4@cref}{{[section][4][16]16.4}{[1][1][]1}}
\xrlabel{sec:16.5}{{16.5}{}{Dopants}{Doc-Start}{}}
\xrlabel{sec:16.5@cref}{{[section][5][16]16.5}{[1][1][]1}}
\xrlabel{eq:16.5-floc}{{16.40}{}{Response of dilute donors}{Doc-Start}{}}
\xrlabel{eq:16.5-floc@cref}{{[equation][40][16]16.40}{[1][1][]1}}
\xrlabel{sec:16.6}{{16.6}{}{Color centers and broken bonds}{Doc-Start}{}}
\xrlabel{sec:16.6@cref}{{[section][6][16]16.6}{[1][1][]1}}
\xrlabel{sec:16.7}{{16.7}{}{Glasses, random media and two-level systems}{Doc-Start}{}}
\xrlabel{sec:16.7@cref}{{[section][7][16]16.7}{[1][1][]1}}
\xrlabel{eq:16.7-f}{{16.60}{}{Dark-matter scattering on two-level systems}{Doc-Start}{}}
\xrlabel{eq:16.7-f@cref}{{[equation][60][16]16.60}{[1][1][]1}}
\xrlabel{eq:16.7-CIR}{{16.66}{}{Vibrational absorption of glasses}{Doc-Start}{}}
\xrlabel{eq:16.7-CIR@cref}{{[equation][66][16]16.66}{[1][1][]1}}
\xrlabel{eq:16.7-AB}{{16.67}{}{Vibrational absorption of glasses}{Doc-Start}{}}
\xrlabel{eq:16.7-AB@cref}{{[equation][67][16]16.67}{[1][1][]1}}
\xrlabel{eq:16.7-Vopt}{{16.70}{}{Random media and the Born limit}{Doc-Start}{}}
\xrlabel{eq:16.7-Vopt@cref}{{[equation][70][16]16.70}{[1][1][]1}}
\xrlabel{eq:16.7-scrit}{{16.71}{}{Random media and the Born limit}{Doc-Start}{}}
\xrlabel{eq:16.7-scrit@cref}{{[equation][71][16]16.71}{[1][1][]1}}
\xrlabel{sec:16.8}{{16.8}{}{Two-level systems at millikelvin}{Doc-Start}{}}
\xrlabel{sec:16.8@cref}{{[section][8][16]16.8}{[1][1][]1}}
\xrlabel{sec:16.9}{{16.9}{}{Material inputs of the novel targets}{Doc-Start}{}}
\xrlabel{sec:16.9@cref}{{[section][9][16]16.9}{[1][1][]1}}
\xrlabel{sec:16.9-limits}{{16.9.2}{}{Theory-limited and data-limited targets}{Doc-Start}{}}
\xrlabel{sec:16.9-limits@cref}{{[subsection][2][16,9]16.9.2}{[1][1][]1}}
\xrlabel{we:16-band}{{16.9}{}{Reach bands from input uncertainties}{Doc-Start}{}}
\xrlabel{we:16-band@cref}{{[workedexample][9][16]16.9}{[1][1][]1}}
\xrlabel{chap:17}{{17}{}{Comparison across channels}{Doc-Start}{}}
\xrlabel{chap:17@cref}{{[chapter][17][]17}{[1][1][]1}}
\xrlabel{sec:17.1}{{17.1}{}{Scattering across channels}{Doc-Start}{}}
\xrlabel{sec:17.1@cref}{{[section][1][17]17.1}{[1][1][]1}}
\xrlabel{sec:17.5}{{17.5}{}{Figures of merit per channel}{Doc-Start}{}}
\xrlabel{sec:17.5@cref}{{[section][5][17]17.5}{[1][1][]1}}
\xrlabel{eq:17.5-partial}{{17.16}{}{Gap floor and valence f-sum}{Doc-Start}{}}
\xrlabel{eq:17.5-partial@cref}{{[equation][16][17]17.16}{[1][1][]1}}
\xrlabel{sec:17.6}{{17.6}{}{Ranking under uncertainty}{Doc-Start}{}}
\xrlabel{sec:17.6@cref}{{[section][6][17]17.6}{[1][1][]1}}
\xrlabel{sec:17.6-correlation}{{17.6.4}{}{Correlated method errors}{Doc-Start}{}}
\xrlabel{sec:17.6-correlation@cref}{{[subsection][4][17,6]17.6.4}{[1][1][]1}}
\xrlabel{chap:18}{{18}{}{Ionization yield}{Doc-Start}{}}
\xrlabel{chap:18@cref}{{[chapter][18][]18}{[1][1][]1}}
\xrlabel{sec:18.1}{{18.1}{}{The carrier cascade}{Doc-Start}{}}
\xrlabel{sec:18.1@cref}{{[section][1][18]18.1}{[1][1][]1}}
\xrlabel{sec:18.2}{{18.2}{}{Mean energy per pair and the Fano factor}{Doc-Start}{}}
\xrlabel{sec:18.2@cref}{{[section][2][18]18.2}{[1][1][]1}}
\xrlabel{eq:18.2-moments}{{18.10}{}{Moments of the heat}{Doc-Start}{}}
\xrlabel{eq:18.2-moments@cref}{{[equation][10][18]18.10}{[1][1][]1}}
\xrlabel{sec:18.3}{{18.3}{}{Seed distributions and the step model}{Doc-Start}{}}
\xrlabel{sec:18.3@cref}{{[section][3][18]18.3}{[1][1][]1}}
\xrlabel{eq:18.3-seed}{{18.14}{}{The seed distribution}{Doc-Start}{}}
\xrlabel{eq:18.3-seed@cref}{{[equation][14][18]18.14}{[1][1][]1}}
\xrlabel{eq:18.3-pcasc}{{18.15}{}{The seed distribution}{Doc-Start}{}}
\xrlabel{eq:18.3-pcasc@cref}{{[equation][15][18]18.15}{[1][1][]1}}
\xrlabel{sec:18.3-fullband}{{18.3.3}{}{The full-band cascade as reference}{Doc-Start}{}}
\xrlabel{sec:18.3-fullband@cref}{{[subsection][3][18,3]18.3.3}{[1][1][]1}}
\xrlabel{eq:18.3-obs}{{18.19}{}{Pair number at the readout}{Doc-Start}{}}
\xrlabel{eq:18.3-obs@cref}{{[equation][19][18]18.19}{[1][1][]1}}
\xrlabel{sec:18.4}{{18.4}{}{Calibrated pair energies}{Doc-Start}{}}
\xrlabel{sec:18.4@cref}{{[section][4][18]18.4}{[1][1][]1}}
\xrlabel{sec:18.4-range}{{18.4.4}{}{The validated range}{Doc-Start}{}}
\xrlabel{sec:18.4-range@cref}{{[subsection][4][18,4]18.4.4}{[1][1][]1}}
\xrlabel{sec:18.5}{{18.5}{}{Capture and excitons}{Doc-Start}{}}
\xrlabel{sec:18.5@cref}{{[section][5][18]18.5}{[1][1][]1}}
\xrlabel{chap:19}{{19}{}{Low-energy nuclear recoils}{Doc-Start}{}}
\xrlabel{chap:19@cref}{{[chapter][19][]19}{[1][1][]1}}
\xrlabel{sec:19.1}{{19.1}{}{Nuclear recoil on the Born--Oppenheimer surface}{Doc-Start}{}}
\xrlabel{sec:19.1@cref}{{[section][1][19]19.1}{[1][1][]1}}
\xrlabel{sec:19.1-handover}{{19.1.1}{}{The hand-over from the nuclear response}{Doc-Start}{}}
\xrlabel{sec:19.1-handover@cref}{{[subsection][1][19,1]19.1.1}{[1][1][]1}}
\xrlabel{sec:19.2}{{19.2}{}{Electronic excitation by the recoil}{Doc-Start}{}}
\xrlabel{sec:19.2@cref}{{[section][2][19]19.2}{[1][1][]1}}
\xrlabel{eq:19.2-kernel}{{19.10}{}{Close collisions and the excitation kernel}{Doc-Start}{}}
\xrlabel{eq:19.2-kernel@cref}{{[equation][10][19]19.10}{[1][1][]1}}
\xrlabel{sec:19.3}{{19.3}{}{Stopping of slow ions}{Doc-Start}{}}
\xrlabel{sec:19.3@cref}{{[section][3][19]19.3}{[1][1][]1}}
\xrlabel{sec:19.4}{{19.4}{}{Defects and stored energy}{Doc-Start}{}}
\xrlabel{sec:19.4@cref}{{[section][4][19]19.4}{[1][1][]1}}
\xrlabel{sec:19.5}{{19.5}{}{The energy balance law}{Doc-Start}{}}
\xrlabel{sec:19.5@cref}{{[section][5][19]19.5}{[1][1][]1}}
\xrlabel{sec:19.6}{{19.6}{}{Delayed release and the low-energy excess}{Doc-Start}{}}
\xrlabel{sec:19.6@cref}{{[section][6][19]19.6}{[1][1][]1}}
\xrlabel{eq:19.6-sumrule}{{19.44}{}{The release-path sum rule}{Doc-Start}{}}
\xrlabel{eq:19.6-sumrule@cref}{{[equation][44][19]19.44}{[1][1][]1}}
\xrlabel{chap:20}{{20}{}{Charge, phonon and quasiparticle readout}{Doc-Start}{}}
\xrlabel{chap:20@cref}{{[chapter][20][]20}{[1][1][]1}}
\xrlabel{sec:20.1}{{20.1}{}{Charge transport and the measured charge}{Doc-Start}{}}
\xrlabel{sec:20.1@cref}{{[section][1][20]20.1}{[1][1][]1}}
\xrlabel{eq:20.1-work}{{20.3}{}{Induced charge and Ramo's theorem}{Doc-Start}{}}
\xrlabel{eq:20.1-work@cref}{{[equation][3][20]20.3}{[1][1][]1}}
\xrlabel{eq:20.1-hecht}{{20.4}{}{Trapping and the Hecht relation}{Doc-Start}{}}
\xrlabel{eq:20.1-hecht@cref}{{[equation][4][20]20.4}{[1][1][]1}}
\xrlabel{eq:20.1-sheet}{{20.9}{}{Space charge of trapped carriers}{Doc-Start}{}}
\xrlabel{eq:20.1-sheet@cref}{{[equation][9][20]20.9}{[1][1][]1}}
\xrlabel{eq:20.1-luke}{{20.10}{}{Luke-phonon emission}{Doc-Start}{}}
\xrlabel{eq:20.1-luke@cref}{{[equation][10][20]20.10}{[1][1][]1}}
\xrlabel{sec:20.2}{{20.2}{}{Athermal phonons in the substrate}{Doc-Start}{}}
\xrlabel{sec:20.2@cref}{{[section][2][20]20.2}{[1][1][]1}}
\xrlabel{sec:20.3}{{20.3}{}{Phonons into films}{Doc-Start}{}}
\xrlabel{sec:20.3@cref}{{[section][3][20]20.3}{[1][1][]1}}
\xrlabel{sec:20.4}{{20.4}{}{Quasiparticle yield in superconductors}{Doc-Start}{}}
\xrlabel{sec:20.4@cref}{{[section][4][20]20.4}{[1][1][]1}}
\xrlabel{sec:20.5}{{20.5}{}{Collection kernels and recoil discrimination}{Doc-Start}{}}
\xrlabel{sec:20.5@cref}{{[section][5][20]20.5}{[1][1][]1}}
\xrlabel{eq:20.5-kernel}{{20.41}{}{Collection kernels}{Doc-Start}{}}
\xrlabel{eq:20.5-kernel@cref}{{[equation][41][20]20.41}{[1][1][]1}}
\xrlabel{eq:21.1-functional}{{21.1}{}{The halo functional of a channel}{Doc-Start}{}}
\xrlabel{eq:21.1-functional@cref}{{[equation][1][21]21.1}{[1][1][]1}}
\xrlabel{eq:21.1-kernel}{{21.2}{}{The halo functional of a channel}{Doc-Start}{}}
\xrlabel{eq:21.1-kernel@cref}{{[equation][2][21]21.2}{[1][1][]1}}
\xrlabel{eq:21.1-gammav}{{21.4}{}{The mean inverse speed}{Doc-Start}{}}
\xrlabel{eq:21.1-gammav@cref}{{[equation][4][21]21.4}{[1][1][]1}}
\xrlabel{eq:21.1-eta2}{{21.6}{}{The mean inverse speed}{Doc-Start}{}}
\xrlabel{eq:21.1-eta2@cref}{{[equation][6][21]21.6}{[1][1][]1}}
\xrlabel{sec:21.2}{{21.2}{}{Halo-independent bounds}{Doc-Start}{}}
\xrlabel{sec:21.2@cref}{{[section][2][21]21.2}{[1][1][]1}}
\xrlabel{eq:21.2-stream}{{21.9}{}{Stream kernels and the linear program}{Doc-Start}{}}
\xrlabel{eq:21.2-stream@cref}{{[equation][9][21]21.9}{[1][1][]1}}
\xrlabel{eq:21.2-column}{{21.11}{}{The bound for a speed-only response}{Doc-Start}{}}
\xrlabel{eq:21.2-column@cref}{{[equation][11][21]21.11}{[1][1][]1}}
\xrlabel{eq:21.2-hib}{{21.12}{}{The bound for a speed-only response}{Doc-Start}{}}
\xrlabel{eq:21.2-hib@cref}{{[equation][12][21]21.12}{[1][1][]1}}
\xrlabel{eq:21.2-multipole}{{21.13}{}{Directional responses in the laboratory}{Doc-Start}{}}
\xrlabel{eq:21.2-multipole@cref}{{[equation][13][21]21.13}{[1][1][]1}}
\xrlabel{sec:21.2-leak}{{21.2.4}{}{Grids, certificates and leakage}{Doc-Start}{}}
\xrlabel{sec:21.2-leak@cref}{{[subsection][4][21,2]21.2.4}{[1][1][]1}}
\xrlabel{eq:21.2-fsleak}{{21.15}{}{Grids, certificates and leakage}{Doc-Start}{}}
\xrlabel{eq:21.2-fsleak@cref}{{[equation][15][21]21.15}{[1][1][]1}}
\xrlabel{sec:21.3-power}{{21.3.3}{}{The edge power law}{Doc-Start}{}}
\xrlabel{sec:21.3-power@cref}{{[subsection][3][21,3]21.3.3}{[1][1][]1}}
\xrlabel{eq:21.3-power}{{21.19}{}{The edge power law}{Doc-Start}{}}
\xrlabel{eq:21.3-power@cref}{{[equation][19][21]21.19}{[1][1][]1}}
\xrlabel{tab:21-edge}{{21.1}{}{Powers of the rate in the dark-matter mass, $\mathcal {R} \propto (\mchi - m_{\mathrm {edge}})^{n}$, at the edges of silicon, for the sharp-cut standard halo ($k_{\mathrm {tail}} = 0$); a halo with tail exponent $k_{\mathrm {tail}}$ adds $k_{\mathrm {tail}}$ to every entry (\cref {eq:21.3-power}). Masses in MeV, with the halo of \cref {tab:1-shm}; $p_{\mathrm {on}}$ is the onset exponent of the direction-averaged response at fixed $|\vect {q}|$; the fourth column says whether $x_{\mathrm {max}}(q)$ has a smooth maximum or a corner at $q^*$ (\cref {eq:21.3-power}). The zero-phonon edges lie above the phonon-assisted onset, where they add a kink to the total rate. The nuclear-recoil row is for a \SI {10}{eV} threshold}{Doc-Start}{}}
\xrlabel{tab:21-edge@cref}{{[table][1][21]21.1}{[1][1][]1}}
\xrlabel{sec:21.4}{{21.4}{}{Theory nuisances in a limit}{Doc-Start}{}}
\xrlabel{sec:21.4@cref}{{[section][4][21]21.4}{[1][1][]1}}
\xrlabel{sec:21.4-like}{{21.4.2}{}{Theory and detector nuisances}{Doc-Start}{}}
\xrlabel{sec:21.4-like@cref}{{[subsection][2][21,4]21.4.2}{[1][1][]1}}
\xrlabel{tab:21-dominant}{{21.2}{}{Dominant theory nuisances of the limit in each channel, read from the nuisance tables of the channel chapters. Sizes are factors or shifts of the rate, and hence, for one bin, of the limit. Basis as in those tables: ``computed'', ``model'', ``est.''\ (parametric estimate), ``spread'' (published calculations at the same order), ``data'' (spread of measured inputs), ``bound'' (only an inequality is known), ``unsized''. The last column gives the section, table or figure where the rows are treated}{Doc-Start}{}}
\xrlabel{tab:21-dominant@cref}{{[table][2][21]21.2}{[1][1][]1}}
\xrlabel{sec:21.5}{{21.5}{}{Combining and recasting limits}{Doc-Start}{}}
\xrlabel{sec:21.5@cref}{{[section][5][21]21.5}{[1][1][]1}}
\xrlabel{sec:21.7}{{21.6}{}{Next targets}{Doc-Start}{}}
\xrlabel{sec:21.7@cref}{{[section][6][21]21.6}{[1][1][]1}}
\xrlabel{tab:21-guide}{{21.4}{}{The guide to next targets. Leaders of the full pool (42 crystals) and of the survivor pool (eleven crystals, \cref {sec:21.7-pools}). Full-pool leader: median with 68\% interval, and probability of ranking first at independent method errors with the value at $r_{\mathrm {mat}} = 1$ in parentheses; the values of the survivor leaders are in \cref {tab:17-rank}. Every probability is conditional on the provenance widths of \cref {tab:17-widths} and on the widths of the pair energies stated in \cref {sec:21.7-pools}; it is not the probability that a crystal is the better target. Missing demonstration: the requirements of \cref {sec:21.7-pools} not documented for the leader. Edges at $\omega _{\mathrm {th}} = \SI {1}{meV}$; $v_{\mathrm {max}} = 2.650\times 10^{-3}$}{Doc-Start}{}}
\xrlabel{tab:21-guide@cref}{{[table][4][21]21.4}{[1][1][]1}}
\xrlabel{chap:22}{{22}{}{Open problems}{Doc-Start}{}}
\xrlabel{chap:22@cref}{{[chapter][22][]22}{[1][1][]1}}
\xrlabel{tab:22-gaps}{{22.1}{}{Exact remaining gaps of the book by chapter, each indexed to the one section that states it. Strongest statement: what that section derives, with its conditions. Missing step: the decisive proof, method, calculation or measurement that would complete it. Limits: \emph {proof}, no number of the book changes. \emph {Bounded}, a number is missing and the section gives a bound or a two-sided bracket. \emph {Estimate}, the section gives an estimate that is not a bound. \emph {Unsized}, the section assigns no size. \emph {Test}, the section computes the number and the missing step checks it. The statements are abbreviated; the section named in the first column gives each in full. Survivors: the eleven crystals with a documented detector record (\cref {sec:21.7}). The calculations that would change a printed result, limit or ranking are ordered in \cref {tab:22-calc}}{Doc-Start}{}}
\xrlabel{tab:22-gaps@cref}{{[table][1][22]22.1}{[1][1][]1}}
\xrlabel{app:D-II}{{E}{}{Long derivations}{Doc-Start}{}}
\xrlabel{app:D-II@cref}{{[appendix][5][2147483647]E}{[1][1][]1}}
\xrlabel{app:D.37}{{E.8}{}{The soft acoustic limit}{Doc-Start}{}}
\xrlabel{app:D.37@cref}{{[subappendix][8][2147483647,5]E.8}{[1][1][]1}}
\xrlabel{eq:D.37-Xi}{{E.30}{}{The comoving vertex}{Doc-Start}{}}
\xrlabel{eq:D.37-Xi@cref}{{[equation][30][2147483647,5]E.30}{[1][1][]1}}
\xrlabel{app:D.43}{{E.13}{}{The cascade generating functional}{Doc-Start}{}}
\xrlabel{app:D.43@cref}{{[subappendix][13][2147483647,5]E.13}{[1][1][]1}}
\xrlabel{eq:D.46-monotone}{{E.53}{}{The step column and the bound}{Doc-Start}{}}
\xrlabel{eq:D.46-monotone@cref}{{[equation][53][2147483647,5]E.53}{[1][1][]1}}
\xrlabel{app:F}{{F}{}{Literature formulas as limits}{Doc-Start}{}}
\xrlabel{app:F@cref}{{[appendix][6][2147483647]F}{[1][1][]1}}
